%%%%%%%%%%%%%%%%%
% This is a two-column CV template based on moderncv.cls
% (v1.0, May 16 2024) written by Genki Ogaki (rockstarogk@gmail.com). Compiles with XeLaTeX.
%
%% It may be distributed and/or modified under the
%% conditions of the LaTeX Project Public License, either version 1.3
%% of this license or (at your option) any later version.
%% The latest version of this license is in
%%    http://www.latex-project.org/lppl.txt
%% and version 1.3 or later is part of all distributions of LaTeX
%% version 2003/12/01 or later.
%%%%%%%%%%%%%%%%

% environment
\documentclass[11pt,a3paper]{moderncv}
\moderncvtheme[green]{banking}
\nopagenumbers{}

\usepackage[T1]{fontenc}
\usepackage{inputenc}
\usepackage[scale=0.9]{geometry}
\usepackage{tabularx}
\usepackage{mathpazo}
\usepackage{hyperref}

\hypersetup{
    colorlinks=true,
    linkcolor=blue,
    urlcolor=blue
}
% macro
\renewcommand*{\labelitemi}{-}

\newcolumntype{L}{>{\raggedright\arraybackslash}X}
\newcolumntype{C}{>{\centering\arraybackslash}X}
\newcolumntype{R}{>{\raggedleft\arraybackslash}X}

\newcommand*{\experienceentry}[5][1.5mm]{
    \subsection{#2} \vspace{-1.5mm}
    \begin{tabularx}{\textwidth}{LR}
        {\itshape #3} & {\itshape #4, #5}
    \end{tabularx}
    \par\addvspace{#1}
}

\newcommand*{\educationentry}[4][0.5mm]{
    \begin{tabularx}{\textwidth}{LR}
        {\bfseries #3} & {\bfseries #4} \\
    \end{tabularx}
    {\itshape #2}
    \par\addvspace{#1}
}

\newcommand*{\scoreentry}[3][2.5mm]{
    {\bfseries #2} \\
    {\itshape #3}
    \par\addvspace{#1}
}

% preamble
\firstname{Abdullah Al Omar}
\familyname{Galib}
\address{Dhaka, Bangladesh}{}

% document
\begin{document}

% Custom fancy header
\begin{center}
    \begin{tikzpicture}[remember picture,overlay]
        \node[rectangle, fill=green!15, anchor=north, minimum width=\paperwidth, minimum height=4.5cm] (box) at (current page.north){};
    \end{tikzpicture}
    \vspace{0.5cm}
    
    {\Huge\bfseries\color{green!50!black} ABDULLAH AL OMAR GALIB}
    
    \vspace{0.3cm}
    {\Large\color{gray} Quantum Computing Researcher \& AI Enthusiast}
    
    \vspace{0.5cm}
    \begin{tabularx}{\textwidth}{*{4}{>{\centering\arraybackslash}X}}
        \color{green!50!black}\faMapMarker & \color{green!50!black}\faEnvelope & \color{green!50!black}\faGithub & \color{green!50!black}\faLinkedin \\
        \small Dhaka, Bangladesh & 
        \small \href{mailto:abdullah.al.omargalib@g.bracu.ac.bd}{abdullah.al.omargalib@g.bracu.ac.bd} & 
        \small \href{https://github.com/ahkatlio}{github.com/ahkatlio} & 
        \small \href{https://www.linkedin.com/in/abdullah-al-omar-galib-30b6b1258/}{LinkedIn Profile}
    \end{tabularx}
    
    \vspace{0.3cm}
    \textcolor{green!50!black}{\rule{0.8\textwidth}{0.5pt}}
\end{center}

\vspace{0.5cm}

% left column
\begin{minipage}[t]{0.62\textwidth}
\section{Experience}
{\small
\experienceentry{Quantum High School Organization}{Co-Founder \& Vice President}{July 2023 - August 2025}{Global}
\begin{itemize}
    \item Spreading quantum education globally through online initiatives, including webinars, workshops, and interactive content.
    \item Mentored 500+ students in quantum computing, creating tailored lesson plans and interactive content.
    \item Organized and led two highly successful quantum education events. The first event was dedicated to teaching quantum computing to beginners, ranging from high school students to graduate students, and was attended by over 300 participants. The second event focused on cutting-edge quantum computing research. Both events received positive feedback and contributed to the growth of quantum education and research.
\end{itemize}
\vspace{1.0mm}

\experienceentry{QWorld}{Research Intern}{July 2023 - August 2023}{Remote}
\begin{itemize}
    \item Researched classical and quantum approaches for Quantum State Tomography (QST).
    \item Implemented and compared different Quantum Machine Learning (QML) approaches for QST, such as supervised and unsupervised learning techniques, to assess their performance and accuracy. Utilized cutting-edge quantum hardware and quantum simulators to conduct QST on multi-qubit networks and complex entangled states.
    \item Achieved an impressive 98\% fidelity in the quantum state reconstruction process, validating the effectiveness and precision of the chosen methodologies.
    \item The research findings and results were documented in a comprehensive research paper, contributing to the scientific community's understanding of quantum state tomography.
\end{itemize}
\vspace{1.0mm}
\experienceentry{iTesseract Technologies Ltd}{Robotics, IoT \& AI Project leader}{Sep 2022 - Nov 2022}{Dhaka, Bangladesh}
\begin{itemize}
    \item Developed and implemented AI software with an outstanding 99\% accuracy rate, significantly reducing the risk of accidents and enhancing overall road safety.
    \item Won the Best Project Director award.
\end{itemize}

}
\vspace{2.0mm}
\section{Projects}
{\small
\newcommand{\project}[3]{
    \item \href{#2}{\textbf{#1}} - #3
}
\begin{itemize}
    \project
        {Bronze Pytket (Quantinuum)}
        {https://gitlab.com/clausia/bronze-pytket}
        {An introductory quantum computing tutorial using Pytket, created in August 2023 with Quantinuum support.}
    \project
        {Quantum Convolutional Neural Networks vs Fully Classical CNNs}
        {https://github.com/ahkatlio/QCNN-vs-CNN-Performance-Analysis}
        {Compared hybrid quantum-classical CNNs with classical CNNs on MNIST, achieving 90.54\% accuracy with the quantum model, outperforming the classical model by 8.57\%.}
    \project
        {Driving Safety (Machine Learning)}
        {https://github.com/ahkatlio/Driving_safety}
        {Created a real-time driving safety application using facial landmark detection and traffic sign recognition, achieving 99\% accuracy.}
    \project
        {Quantum Computing Notebook for Beginners}
        {https://github.com/ahkatlio/QHSO_Basics_of_Qiskit}
        {Authored a beginner-friendly Jupyter notebook using Qiskit to introduce fundamental quantum computing concepts.}
    \project
        {Fruit Detection and Classification}
        {https://github.com/ahkatlio/Fruit-and-vegetable-image-recognition}
        {Created an automated image recognition system using YOLOv8 to detect and classify 36 classes of fruits and vegetables with 91.68\% mean Average Precision (mAP). Employing computer vision and machine learning techniques, it is optimized for real-time use in agricultural monitoring, quality control processes, inventory management, and educational tools to improve efficiency in produce handling.}
    \project
        {Nighthawk: AI-Powered Modular Security Tool Orchestrator}
        {https://github.com/ahkatlio/Nighthawk}
        {Built an AI-powered system to orchestrate modular security tools, facilitating integrated management of cybersecurity components for threat detection and response in dynamic environments.}
    \project
        {Quantum Tunneling Simulation with QuTiP and Manim}
        {https://github.com/ahkatlio/quantum-tunneling-sim-with-qutip-and-manim}
        {Developed a simulation to model and visualize quantum tunneling phenomena, utilizing QuTiP for solving quantum dynamics equations and Manim for creating animated representations of wave functions and particle behaviors.}
    \project
        {C\_Flex: A Collection of Whimsical Python Tools}
        {https://github.com/ahkatlio/C_Flex}
        {Assembled a diverse set of Python utilities and experiments as a creative coding playground, including a terminal-based Markdown previewer with navigation and rendering features, AI tools for math solving and chatting, games like Conway's Game of Life, media downloaders, and a WiFi scanner. Employed libraries such as Rich and Curses for UI, emphasizing modular development, experimentation, and cross-platform compatibility via installation scripts.}

\end{itemize}
}
\end{minipage}
\hfill
% right column
\begin{minipage}[t]{0.35\textwidth}
{\small
\nocite{*}
\bibliographystyle{plainurl} % or IEEEtran
\bibliography{publications}
}


\section{Skills}
{\small
\textbf{Quantum Computing:} Qiskit, Pytket, criq\\
\textbf{AI \& Machine Learning:} TensorFlow, Keras \\
\textbf{Software Development:} Unreal Engine 5, Python \\
\textbf{Community \& Education: }Webinars, Workshops, Lesson Planning, Mentoring \\
\textbf{Project Management:} Project Leadership, Team Collaboration, Time Management \\
\textbf{Soft Skills:} Leadership, Communication, Collaboration, Problem-Solving
}
\section{Education}
{\small
\educationentry{Bachelor of Science - Theoretical Physics}{BRAC University}{Dhaka, Bangladesh}
Currently pursuing
\par
\vspace{3.0mm}
\educationentry{Higher Secondary School Certificate (HSC)}{St. Joseph Higher Secondary School}{Dhaka, Bangladesh}
Result: GPA 5.00
}
\section{Languages}
\begin{itemize}
    \item Bangla (Native)
    \item English (Advanced)
\end{itemize}


\section{Awards}
{\small
\begin{itemize}
    \item \textbf{Best Project and Best Presentation awards}, Issued by QWorld, Aug 2023
    \item Successfully participated in \textbf{unitaryHACK2025} on May 28 - June 11, 2025, and contributed to the quantum open-source ecosystem.
    \item Discover more about my accolades and achievements by following this \href{https://drive.google.com/drive/folders/1n1txPoDAXv5QNDxpmu_m7q72QB5I_qZ_}{LINK}.
\end{itemize}
}
\end{minipage}

\end{document}
